\section{Public Research Contract}

The paper has two contracts. The theoretical contract is that every option position is
defined by a mark, a payoff convention, a NAV normalization, a Greek exposure map, and a
residual risk term. The empirical contract is that every number in the result section is
regenerated from files inside the repository. This is a mechanical constraint, but it has
a substantive purpose: it prevents the article from drifting away from the code and data
that support it.

The investable objects are not strategies such as covered calls, protective puts, or
spreads. They are individual listed equity options grouped into rolling buckets only for
empirical continuity. The optimizer sees arbitrary call and put assets, signed position
weights, expected returns, a Greek-induced covariance matrix, and constraints. This is
important because the paper is not a defense of any named option strategy. It is an
allocation model for an arbitrary option universe.

The accounting unit is portfolio NAV. If \(q_i=0.10\), the portfolio pays ten cents of
NAV to buy option \(i\); if \(q_i=-0.10\), it sells an option exposure with ten cents of
initial premium magnitude. Let beginning NAV be \(W_t\), option mark be \(C_{i,t}\), and
contract count be \(N_{i,t}=q_{i,t}W_t/C_{i,t}\). The next-period contribution to NAV is
\[
  \frac{N_{i,t}(C_{i,t+1}-C_{i,t})}{W_t}
  =
  q_{i,t}\frac{C_{i,t+1}-C_{i,t}}{C_{i,t}} .
\]
Thus the option premium is not ignored. It is the denominator of the option return and
the scale of the NAV weight. A long option that expires worthless contributes \(-q_i\) to
NAV, exactly as a paid premium should. The unspent or collateral cash is an accounting
numeraire, not an optimized risk asset. This differs from a stock-cash allocation problem:
the paper does not search for an optimal cash sleeve, risk-free lending amount, or
leverage target. It searches for option premium exposures under gross and risk budgets.

The timing rule is also part of the contract. Expected returns, underlying covariance,
volatility-shock covariance, and residual covariance are estimated from the training window
only. The test-window portfolio returns are then computed from held-out option-bucket
returns. The selected test option is chosen at the decision snapshot; the later
split-adjusted listed-expiry payoff is used only to realize the return. Equity benchmarks
and factor regressions use the same train/test split. No result uses future test returns
to set the Greek Markowitz weights.

The claim hierarchy is deliberately conservative:
\begin{itemize}[leftmargin=2em]
  \item mathematical claims require a displayed derivation or proof;
  \item source claims require a cited paper, book, or primary data document;
  \item data claims require a generated table, figure, CSV, or JSON summary;
  \item production claims would require executable fills, costs, margin, assignment, and
  broker-level parity checks, and are not made here.
\end{itemize}

This distinction controls the strength of the conclusion. The paper can claim that the
option-only Markowitz covariance object is mathematically coherent. It can claim that the
local empirical implementation is point-in-time auditable and produces informative risk
diagnostics on the available OPRA-derived panel. It cannot claim live trading alpha,
capacity, or production tradability.
